\begin{figure}[tbp]
\centering
\begin{tikzpicture}[gclplate,x=1cm,y=1cm]
  \node[gcllabel,draw=GCLBlue!75!black,align=center,text width=2.8cm,minimum height=1.0cm] (a) at (-4.15,0) {synchronous\\send+receive};
  \node[gcllabel,draw=GCLWarm!75!black,align=center,text width=2.8cm,minimum height=1.0cm] (b) at (0,0) {different semantics};
  \node[gcllabel,draw=GCLGreen!75!black,align=center,text width=2.8cm,minimum height=1.0cm] (c) at (4.15,0) {asynchronous\\enqueue then dequeue};
  \draw[gclbluearrow] (a.east) -- (b.west);
  \draw[gclwarmarrow] (b.east) -- (c.west);
\end{tikzpicture}
\caption{Plate 15. Handshake versus queue. A rendezvous and a buffered exchange have different intermediate states. No global trace-equivalence theorem is asserted.}
\label{plate:protocol-15}
\end{figure}
